\chapter{LX-020 - DOGFOOD Locale Demo Readiness}\label{record-196}

\noindent{\footnotesize\color{BookMuted}\path{docs/design/LX-020-dogfood-locale-demo-readiness.md}\par}\smallskip

\begin{RecordMeta}
id: LX-020
title: DOGFOOD Locale Demo Readiness
status: active
lane: asap
priority: high
github_issue: 340
legend: LX

\end{RecordMeta}

Legend: \href{https://github.com/flyingrobots/bijou/blob/1853ad902f6bf219326409e7f743a5a1d6517dbc/docs/legends/LX-localization-and-bidirectionality.md}{LX - Localization and Bidirectionality}

\section{Decision Summary}\label{record-196-decision-summary}

DOGFOOD should stop presenting non-English locale selection as if the whole app is fully translated. For \texttt{v\allowbreak{}7\allowbreak{}.\allowbreak{}2\allowbreak{}.\allowbreak{}0}, the honest release-demo contract is:

\begin{itemize}
\tightlist
\item
  core app chrome and release-demo navigation do not show missing localization placeholders
\item
  English-source documentation is labeled as such when selected locale content is intentionally unavailable
\item
  localization debt remains visible to maintainers without leaking broken UI into the user-facing demo path
\end{itemize}

\section{Sponsored Human}\label{record-196-sponsored-human}

A release-video user needs DOGFOOD to survive a non-English locale switch without looking broken. Missing localization markers, accidental English chrome, and unclear partial-translation behavior make the shipped story feel unfinished.

\section{Sponsored Agent}\label{record-196-sponsored-agent}

An agent needs deterministic localization witnesses that distinguish three states:

\begin{itemize}
\tightlist
\item
  selected-locale translation exists and is rendered
\item
  selected-locale translation is intentionally unavailable and the app presents an English-source notice
\item
  selected-locale translation is missing by accident and should fail the gate
\end{itemize}

\section{Hill}\label{record-196-hill}

After this cycle, DOGFOOD's non-English demo path is truthful enough for a release video: switching locales does not expose missing-placeholder strings in major app chrome, and documentation surfaces clearly explain when they remain English-source content.

\section{Problem}\label{record-196-problem}

Issue \#340 came from the \texttt{v\allowbreak{}7\allowbreak{}.\allowbreak{}1\allowbreak{}.\allowbreak{}0} release-video rehearsal. DOGFOOD can switch to non-English locales, but much of the app still renders English copy or missing localization markers. Existing i18n gates prove debt is counted, not that the presentation path is demo-ready.

The current debt ratchet is useful, but it is not sufficient for user-facing release claims:

\begin{CodeBlock}
dogfood-i18n-debt: ok (2772 raw strings; baseline 2772; 78 missing Markdown localizations; baseline 78)
\end{CodeBlock}

That output tells maintainers the known debt did not grow. It does not tell a viewer why non-English DOGFOOD still shows English documentation.

\section{Scope}\label{record-196-scope}

\begin{itemize}
\tightlist
\item
  Audit the launch, navigation, settings, status/footer, and major tab chrome for the supported non-English DOGFOOD locales.
\item
  Add deterministic coverage that fails when the release-demo path renders visible missing-localization placeholders.
\item
  Add or update DOGFOOD copy so intentionally English-source documentation is labeled clearly when the selected locale lacks translated Markdown.
\item
  Keep the existing debt ratchet intact; this cycle strengthens the presentation contract around it.
\item
  Update release-facing docs/changelog so \texttt{v\allowbreak{}7\allowbreak{}.\allowbreak{}2\allowbreak{}.\allowbreak{}0} describes the localization posture honestly.
\end{itemize}

\section{Non-Goals}\label{record-196-non-goals}

\begin{itemize}
\tightlist
\item
  Do not translate the full DOGFOOD documentation corpus in this cycle.
\item
  Do not build a localization workbench or dashboard.
\item
  Do not change the CSV/XLSX localization workflow.
\item
  Do not hide known debt from maintainer tooling.
\item
  Do not broaden \texttt{v\allowbreak{}7\allowbreak{}.\allowbreak{}2\allowbreak{}.\allowbreak{}0} into the V8/V9 localization operations surface.
\end{itemize}

\section{User Experience Contract}\label{record-196-user-experience-contract}

For non-English locales, DOGFOOD may show English documentation only if the UI makes that fallback explicit. The user should not have to infer whether English copy is a bug, an untranslated source document, or a missing catalog entry.

Major app chrome should be localized or explicitly excluded from the release demo path. Missing-placeholder strings are not acceptable on the demo path.

\section{Lower Modes}\label{record-196-lower-modes}

\subsection{Static Mode}\label{record-196-static-mode}

Static DOGFOOD render helpers should be able to render a selected-locale path and expose whether missing placeholders appear in visible output.

\subsection{Pipe Mode}\label{record-196-pipe-mode}

Pipe/smoke output should remain deterministic and suitable for release evidence. Locale-specific output checks must avoid terminal-profile assumptions.

\subsection{Accessible Mode}\label{record-196-accessible-mode}

English-source notices should be text, not color-only styling. Screen readers and terminal readers should receive the same partial-localization status.

\section{Localization / Directionality Posture}\label{record-196-localization-directionality-posture}

This cycle treats localization status as product truth. The app may still carry known translation debt, but the release-demo path must make that debt explicit instead of looking accidentally broken.

No new right-to-left layout contract is introduced. Direction metadata should continue to flow from the selected locale through the existing localization port/runtime path.

\section{Agent Inspectability / Explainability Posture}\label{record-196-agent-inspectability-explainability-posture}

Agents should be able to inspect:

\begin{itemize}
\tightlist
\item
  which locale was selected
\item
  whether a rendered string came from selected-locale content, fallback content, or an English-source documentation notice
\item
  whether missing localization markers appeared in the tested output
\end{itemize}

The proof should be executable through focused tests and existing DOGFOOD i18n commands, not through screenshots alone.

\section{Linked Invariants}\label{record-196-linked-invariants}

\begin{itemize}
\tightlist
\item
  DOGFOOD visible release claims must be true in the app.
\item
  Known localization debt must stay counted by the ratchet.
\item
  Missing localization placeholders must not appear in the release-demo path.
\item
  English-source documentation fallback must be explicit to users.
\end{itemize}

\section{Implementation Plan}\label{record-196-implementation-plan}

\begin{enumerate}
\def\labelenumi{\arabic{enumi}.}
\tightlist
\item
  Add focused failing tests that render representative non-English DOGFOOD paths and assert no visible missing-localization placeholders are present.
\item
  Identify the release-demo chrome surfaces that should be localized now.
\item
  Add an English-source documentation notice for intentionally untranslated Markdown/documentation surfaces.
\item
  Update catalog entries or DOGFOOD copy for the scoped chrome surfaces.
\item
  Preserve the existing \texttt{d\allowbreak{}o\allowbreak{}g\allowbreak{}f\allowbreak{}o\allowbreak{}o\allowbreak{}d\allowbreak{}:\allowbreak{}i\allowbreak{}1\allowbreak{}8\allowbreak{}n\allowbreak{}:\allowbreak{}d\allowbreak{}e\allowbreak{}b\allowbreak{}t}, \texttt{d\allowbreak{}o\allowbreak{}g\allowbreak{}f\allowbreak{}o\allowbreak{}o\allowbreak{}d\allowbreak{}:\allowbreak{}i\allowbreak{}1\allowbreak{}8\allowbreak{}n\allowbreak{}:\allowbreak{}c\allowbreak{}h\allowbreak{}e\allowbreak{}c\allowbreak{}k}, and \texttt{d\allowbreak{}o\allowbreak{}g\allowbreak{}f\allowbreak{}o\allowbreak{}o\allowbreak{}d\allowbreak{}:\allowbreak{}i\allowbreak{}1\allowbreak{}8\allowbreak{}n\allowbreak{}:\allowbreak{}c\allowbreak{}o\allowbreak{}m\allowbreak{}p\allowbreak{}l\allowbreak{}e\allowbreak{}t\allowbreak{}e} gates.
\item
  Update changelog and any release-facing wording needed to avoid overclaiming full DOGFOOD translation.
\end{enumerate}

\section{Tests To Write First}\label{record-196-tests-to-write-first}

\begin{itemize}
\tightlist
\item
  A non-English DOGFOOD launch/render path fails if visible missing localization markers appear.
\item
  A selected non-English documentation path renders an explicit English-source notice when translated Markdown is unavailable.
\item
  Locale switching keeps major app chrome free of accidental missing markers.
\item
  Existing debt-ratchet behavior still reports known raw-string and Markdown localization debt.
\end{itemize}

\section{Acceptance Criteria}\label{record-196-acceptance-criteria}

\begin{itemize}
\tightlist
\item
  Non-English release-demo path has deterministic missing-placeholder coverage.
\item
  Major DOGFOOD chrome on that path is localized or explicitly excluded from the release-demo claim.
\item
  English-source documentation fallback is visible and textual.
\item
  Existing i18n debt gates remain green.
\item
  \texttt{d\allowbreak{}o\allowbreak{}c\allowbreak{}s\allowbreak{}/\allowbreak{}C\allowbreak{}H\allowbreak{}A\allowbreak{}N\allowbreak{}G\allowbreak{}E\allowbreak{}L\allowbreak{}O\allowbreak{}G\allowbreak{}.\allowbreak{}m\allowbreak{}d} records the demo-readiness fix.
\item
  Issue \#340 is closed by the implementation PR.
\end{itemize}

\section{Retrospective}\label{record-196-retrospective}

To be completed when the PR lands.

